% signed_border_mean.tex -- round 444, the signed mean of
% q^dagger as a triple sum, with a circularity gate
% (last analytic object of the RH path after r442).
% ZIRKULAER / TRIPLE_SUM_EXACT / POLE_NOT_CARRIER /
% DIAGONAL_OVERFLOWS / OFFDIAG_LOADBEARING /
% DEAD_CHI_POLE_OVERSHOOT.
% Builds with pdflatex.
% Companion machine check: verify_signed_border.py.
% Discovery: experiments/tfpt-discovery/signed_border_mean_probe.py.
% NO RH CLAIM.  Research documentation, not a theorem of RH.
\documentclass[11pt]{article}

\usepackage[a4paper,margin=2.7cm]{geometry}
\usepackage{amsmath,amssymb,amsthm}
\usepackage{microtype}
\usepackage{booktabs}
\usepackage{xcolor}
\usepackage[colorlinks=true,linkcolor=blue!50!black,urlcolor=blue!50!black]{hyperref}

\newtheorem{lemma}{Lemma}
\newtheorem{prop}{Proposition}
\theoremstyle{remark}
\newtheorem{remark}{Remark}

\newcommand{\Rd}{R^{\dagger}}
\newcommand{\qd}{q^{\dagger}}
\newcommand{\kapd}{\kappa^{\dagger}}
\newcommand{\indm}{\operatorname{ind}_{-}}
\newcommand{\half}{\tfrac12}

\title{\bfseries The signed border mean\\[2mm]
\large triple sum exact; pole not the carrier;\\
the bound sits at the circularity gate}
\author{Stefan Hamann}
\date{August 30, 2026}

\begin{document}
\maketitle

\begin{abstract}\noindent
r442 reduced the remaining RH path to a positive frequency
of $\qd<1$ on dyadic $\zeta$-anchors, with
$\qd=B_w^{-1}\,b^{T}(I-C)^{-1}b$.
This note writes the block mean as an explicit triple sum
over windows and signed border atoms, splits diagonal from
off-diagonal and the almost-mode pole from the regular
resolvent, and attempts a classical mean-value bound.
The identities are theorems of finite linear algebra.
The bound is \emph{circular}: the load-bearing object is
signed off-diagonal cancellation in the \emph{regular}
kernel, of strength matching a diagonal that already
exceeds $40$ at $k=9$, against a window-dependent
OP-resolvent kernel that is not a classical
$k(\log n/n')$.  That is the second precise circularity
verdict of the path (after r399).  No RH claim.
\end{abstract}

\medskip\noindent
\textbf{Status.}\enspace
Lemmas~\ref{lem:triple}--\ref{lem:pole} are \textbf{SATZ}.
Propositions~\ref{prop:diag}--\ref{prop:chi} are
machine-checked census or named refutations.
The mean bound is \emph{circular}, not open-untested:
the missing signed strength is named in
Section~\ref{sec:bound}.
Ausgang
\texttt{ZIRKULAER / TRIPLE\_SUM\_EXACT / POLE\_NOT\_CARRIER / DIAGONAL\_OVERFLOWS / OFFDIAG\_LOADBEARING / DEAD\_CHI\_POLE\_OVERSHOOT}.
No $L^{*}$ claim.  No $R^{\dagger}$ claim.  No RH claim.

\section{The triple sum}\label{sec:triple}

r420/r424 write $b$ as the $\mu$-OP coefficient of the
signed border measure $\sigma$ on atoms $(t_a,\sigma_a)$:
$b=\sum_a\sigma_a\,\phi(t_a)$.
Woodbury gives $\qd=B_w^{-1}\,b^{T}(I-C)^{-1}b$ with
$C=B^{T}B$.  Expanding the quadratic form:

\begin{lemma}[Triple sum]\label{lem:triple}
On each window $W_k$
\[
\qd(W_k)
=\sum_{a,a'}
\frac{\sigma_{k,a}\sigma_{k,a'}}{B_{w,k}}
\bigl\langle
\phi_k(t_a),\,(I-C_k)^{-1}\phi_k(t_{a'})
\bigr\rangle,
\]
hence
\[
\frac1K\sum_{k=1}^K\qd(W_k)
=\frac1K\sum_{k}\sum_{a,a'}
\frac{\sigma_{k,a}\sigma_{k,a'}}{B_{w,k}}
K_k(a,a'),
\]
where $K_k(a,a')=\langle\phi_k(t_a),(I-C_k)^{-1}\phi_k(t_{a'})\rangle$.
Over $\mathbb{Q}$: two atoms $\phi=(1,0)$ and $(1,1)$
against $C=\mathrm{diag}(1/2,0)$ give $q=9=5+4$.
On w9 the reconstruction $b=\Phi\sigma$ holds at
$1.4\cdot 10^{-16}$ and matches $\qd$ at $10^{-15}$.
\end{lemma}

Selected anchors $a_k=2^k$ are $\zeta$-windows: the
$X$-atoms include primes in dyadic ranges (von~Mangoldt
data), the $Y$-atoms are holes with opposite sign.

\section{Diagonal versus off-diagonal}\label{sec:diag}

Write $\qd=q_{\mathrm{diag}}+q_{\mathrm{off}}$ with
$q_{\mathrm{diag}}=B_w^{-1}\sum_a\sigma_a^{2}\,K(a,a)$
the unsigned skeleton (signs drop).

\begin{prop}[Diagonal overflows]\label{prop:diag}
Selected $k=3,4,5,6,7,9$:
\[
\begin{array}{crrrr}
k&q_{\mathrm{diag}}&q_{\mathrm{off}}&\qd&q_{\mathrm{diag}}<1?\\
3&0.515&+0.413&0.928&\mathrm{yes}\\
4&0.482&+0.451&0.933&\mathrm{yes}\\
5&0.892&-0.001&0.891&\mathrm{yes}\\
6&0.835&+0.097&0.931&\mathrm{yes}\\
7&2.272&-1.327&0.945&\mathrm{no}\\
9&40.777&-39.815&0.962&\mathrm{no}\\
\end{array}
\]
Mean $q_{\mathrm{diag}}=7.629>1$.
From $k=7$ the signed off-diagonal is load-bearing:
without it, $\qd>1$ already at $k=7$.
\end{prop}

A positive diagonal majorant therefore cannot prove
$\frac1K\sum\qd<1$.  This is the same family as the r442
unsigned envelopes ($|b|=38.5$ on w9, $|\sigma|=2.10$ at
$k=5$ and $87$ at $k=9$), now seen inside the signed
expansion: even $\sigma^{2}$ on the kernel diagonal
overflows.  The arithmetic is the off-diagonal
cancellation $q_{\mathrm{off}}\approx -q_{\mathrm{diag}}+O(1)$.

The pairwise phases $\log t_a-\log t_{a'}$ are the
classical bilinear differences.  On an $80$-atom
$X$-subsample of w9,
$\mathrm{corr}(K_{\mathrm{off}},\cos\Delta\log)=0.026$:
the kernel is \emph{not} a cosine of the log-ratio.

\section{Pole versus regular resolvent}\label{sec:pole}

Let $\lambda_{\max}(C_k)<1$ (true on every measured living
and dead window: $n_{\ge 1}=0$) with unit eigenvector
$\psi_{\max}$.  As in r436,
\[
(I-C)^{-1}
=\frac{\psi_{\max}\psi_{\max}^{T}}{1-\lambda_{\max}}
+R_{\mathrm{reg}}.
\]
Then $\qd=q_{\mathrm{pole}}+q_{\mathrm{reg}}$ with
$q_{\mathrm{pole}}=B_w^{-1}\lvert\langle b,\psi_{\max}\rangle\rvert^{2}/(1-\lambda_{\max})$.

\begin{lemma}[Spectral split]\label{lem:pole}
Over $\mathbb{Q}$: $C=\mathrm{diag}(1/2,1/4)$, $b=(1,1)$
give $q=10/3=2+4/3$.
On every selected window the split residual is
$\le 10^{-12}$.
\end{lemma}

\begin{prop}[Pole is not the carrier]\label{prop:pole}
Selected pole shares:
$0.0003,\,0.0007,\,0.0031,\,0.0009,\,0.0005,\,0.0001$.
Mean $q_{\mathrm{pole}}=8.6\cdot 10^{-4}$.
The regular part \emph{is} $\qd$ to $10^{-3}$, and follows
the same rise toward $1$ ($0.928\to 0.962$).
The $C$-gap $1-\lambda_{\max}$ shrinks as the r440 collar
($2.99\cdot 10^{-3}\to 1.91\cdot 10^{-8}$), but
$\langle b,\psi_{\max}\rangle$ shrinks faster: the
almost-mode is nearly orthogonal to the border vector on
living windows.
A near-$1$ cluster grows ($65$ modes with gap $<10^{-2}$ at
$k=9$) yet carries only $0.8$~percent of $\qd$.
The living mean is a \emph{bulk regular} quadratic form,
not a pole-weighted prime sum.
\end{prop}

The hoped-for r436-style reduction (pole carries growth,
remainder block-stable $\Rightarrow$ the mean is a
pole-weighted $\Lambda$-sum) is therefore refuted as a
strategy for the living mean.  The remainder is not
block-stable: it \emph{is} the trending $\qd$.

\section{Dead $\chi$: pole overshoot}\label{sec:chi}

\begin{prop}[Death certificate]\label{prop:chi}
On all six dead $\chi$ windows,
$\qd>1$, $q_{\mathrm{reg}}<1$, and
$q_{\mathrm{pole}}>1-q_{\mathrm{reg}}$.
CHI$_3$-$15$: $\qd=1.023$, $q_{\mathrm{pole}}=0.038$,
$q_{\mathrm{reg}}=0.985$.
CHI$_3$-$39$ has the largest pole ($0.095$).
Living CHI$_3$-$9$: $q_{\mathrm{pole}}=0.003$.
Dropping the pole misclassifies every dead window as
living.
\end{prop}

So the almost-mode \emph{does} distinguish live from dead,
but as a death spike, not as the carrier of the living
mean.  On MAIN (core-$42$) the risk is the opposite:
$q_{\mathrm{reg}}$ itself sits at $0.9985$ (kz$16$), with
mean top-mode $4.5\cdot 10^{-4}$.  Two mechanisms, only
one seen as death in the census.

\section{Bound inventory and the gate}\label{sec:bound}

Each brick is typed.  None closes.

\begin{itemize}
\item \textbf{Diagonal majorant} (unconditional).
REFUTED: mean $q_{\mathrm{diag}}=7.63>1$, and
$q_{\mathrm{diag}}>1$ pointwise from $k=7$.
\item \textbf{Pole-weighted prime sum} (would be PNT on
the main term of $\langle\sigma,\psi_{\max}\rangle$,
RH-near on the oscillation).
REFUTED as a reduction: pole share $<0.4$~percent on
Selected.  PNT is idle.
\item \textbf{Montgomery--Vaughan MVT} on
$\sum_{a\neq a'}\sigma_a\sigma_{a'}K_k(a,a')$
(unconditional).
The kernel $K_k$ depends on the window through $C_k$ and
the OP of $W_k$.  It is not a translation-invariant
$k(\log n/n')$ ($\mathrm{corr}=0.026$ vs.\ cosine).
MVT does not apply in the form that would bound a
classical bilinear prime form.  Even if forced, MVT
gives a mean-square envelope of size $\asymp q_{\mathrm{diag}}$,
not cancellation to relative error $1/40$.
\item \textbf{Cauchy--Schwarz} between block and atoms
(unconditional).
Loses the sign of $q_{\mathrm{off}}$.  At $k=9$,
$|q_{\mathrm{off}}|/q_{\mathrm{diag}}=0.976$
(near-perfect cancellation).  CS cannot see it.
\item \textbf{Gallagher large sieve} (unconditional).
Wants a Fourier multiplier in $t=\log$.  The kernel is
not one.  Structurally inapplicable.
\item \textbf{RH-quality bilinear estimate} on
$\sum\Lambda(n)\Lambda(n')\,R^{\mathrm{reg}}_k(n,n')$.
This \emph{would} give the missing cancellation.
It is the circularity gate: an error term of that
quality on signed von-Mangoldt correlations is RH-near
(the zeros of $\zeta$ \emph{are} those correlations).
Forbidden as a mean-value carrier, as in r399.
\end{itemize}

\paragraph{The missing signed strength.}
A uniform bound
\[
\Biggl|
\frac1K\sum_k\sum_{a\neq a'}
\frac{\sigma_{k,a}\sigma_{k,a'}}{B_{w,k}}
R^{\mathrm{reg}}_k(a,a')
\Biggr|
\]
of strength matching the diagonal growth
($q_{\mathrm{diag}}$ already $40$ at $k=9$), with
$R^{\mathrm{reg}}_k=(I-C_k)^{-1}$ minus the rank-one
almost-mode, a window-dependent OP-resolvent kernel on
dyadic $\zeta$-anchors.
This is the second precise circularity verdict of the
path (r399: Weyl-energy decay at CD-transfer scale was
RH-near \emph{and} independently refuted by a growing
census; here the census is living, the bound itself is
circular).

r430 \texttt{exists\_index\_zero\_of\_block\_mean\_lt\_one}
is already proved and is not in question.  The analytic
input it needs is not supplied by any unconditional
brick above.

\section{Kills}\label{sec:kills}

\begin{itemize}
\item \textbf{Unsigned regression.}
$|b|$-envelope w9 $=38.49$; $|\sigma|$ at $k=5,9$ is
$2.10$ and $87.1$ (r442 reproduced).
\item \textbf{Drop-pole mutant.}
$q_{\mathrm{reg}}<1$ on $6/6$ dead $\chi$: the mutant
declares false life.  On living windows drop-pole is
invisible (pole share $\sim 0$).
\item \textbf{Wrong block normalisation.}
$q\cdot B_w=7.82$ on w9, not a $<1$ object.
\item \textbf{Soft means (r440).}
$\operatorname{tr}G^{\dagger}$ is $50.9$ living vs.\ $53.0$
dead; $M_2=32.96$ vs.\ $39.01$.  Both $\gg$ threshold.
\item \textbf{Scramble.}
Unaugmented $\mathrm{nneg}(R)=21$.
\end{itemize}

\section*{Claim boundary}

Nothing in this note is evidence for or against the
Riemann hypothesis, in either direction.
Lemmas~\ref{lem:triple}--\ref{lem:pole} are finite linear
algebra.  The mean-value bound that would feed r430 sits
at the circularity gate named in Section~\ref{sec:bound}.

\begin{thebibliography}{9}
\bibitem{r399} Round 399, \texttt{source\_weyl\_energy}:
  first circularity verdict (Weyl energy).
\bibitem{r436} Round 436, \texttt{p2\_determinant}:
  pole / regular template.
\bibitem{r440} Round 440, \texttt{mean\_tau\_index}:
  collar gap; soft means blind.
\bibitem{r442} Round 442, \texttt{block\_mean}:
  $\kapd=1\{\qd>1\}$; unsigned majorant refuted.
\bibitem{r430} \texttt{exists\_index\_zero\_of\_block\_mean\_lt\_one}.
\end{thebibliography}

\end{document}
